% ─────────────────────────────────────────────────────────────────────────────
%  Capítulo 1 – Introducción
% ─────────────────────────────────────────────────────────────────────────────
\chapter{Introducción}
\label{chap:intro}

\section{Motivación y objetivos}

El diseño e implementación de un compilador desde cero constituye uno de los
proyectos más completos en el ámbito de la Ingeniería del Software, pues pone
en juego simultáneamente teoría de lenguajes formales, estructuras de datos
avanzadas, generación de código máquina y, en este caso, lógica matemática
aplicada a la verificación de programas.

El proyecto \textbf{JavaSubset Compiler} tiene como objetivo construir un
compilador completo para un subconjunto estáticamente tipado de Java con destino
a código nativo x86-64. El proceso se organiza en cinco fases o \emph{milestones}:

\begin{enumerate}[label=\textbf{Fase \arabic*.}]
  \item \textbf{Frontend}: análisis léxico y sintáctico, construcción del AST.
  \item \textbf{Análisis semántico}: resolución de nombres, comprobación de tipos.
  \item \textbf{Representación intermedia}: traducción a código de tres direcciones.
  \item \textbf{Backend x86-64}: generación de ensamblador, asignación de registros.
  \item \textbf{Verificación formal}: Lógica de Hoare, cálculo WP, integración con Z3.
\end{enumerate}

\section{Subconjunto del lenguaje}

El lenguaje fuente es un subconjunto tipado de Java con las siguientes
características:

\begin{table}[H]
  \centering
  \caption{Características del subconjunto Java soportado}
  \label{tab:lenguaje}
  \begin{tabularx}{\textwidth}{lX}
    \toprule
    \textbf{Categoría}   & \textbf{Construcciones soportadas} \\
    \midrule
    Tipos primitivos     & \texttt{int}, \texttt{double}, \texttt{boolean} \\
    Declaraciones        & Variables locales, constantes \texttt{final} \\
    Operadores           & $+$, $-$, $*$, $/$ , $\%$, comparación, lógicos \\
    Control de flujo     & \texttt{if}/\texttt{else}, \texttt{while} \\
    Métodos              & Estáticos sin parámetros, \texttt{return} \\
    E/S estándar         & \texttt{System.out.print/println}, \texttt{Scanner} \\
    Anotaciones formales & \texttt{@pre}, \texttt{@post}, \texttt{@invariant}, \texttt{@variant} \\
    \bottomrule
  \end{tabularx}
\end{table}

\section{Herramientas utilizadas}

\begin{description}
  \item[\textbf{Flex 2.6}] Generador de analizadores léxicos a partir de
    expresiones regulares. Produce la función \texttt{yylex()} en C.
  \item[\textbf{GNU Bison 3.x}] Generador de analizadores sintácticos LALR(1).
    Produce el autómata de análisis y las acciones semánticas en C.
  \item[\textbf{GCC / G++ 11}] Compilador del código C/C++ generado.
    Se usa también para ensamblar y enlazar el código x86-64 producido.
  \item[\textbf{Z3 4.x}] Solver SMT de Microsoft Research. Recibe fórmulas
    en formato SMT-Lib2 y determina su satisfacibilidad (SAT/UNSAT).
\end{description}

\section{Estructura del documento}

El documento sigue la estructura del pipeline del compilador:
el capítulo~\ref{chap:arquitectura} presenta la vista global del sistema;
los capítulos~\ref{chap:frontend}–\ref{chap:backend} describen cada fase
de compilación; el capítulo~\ref{chap:optimizacion} detalla las
optimizaciones sobre la RI; el capítulo~\ref{chap:registro} explica la
asignación de registros; el capítulo~\ref{chap:verificacion} desarrolla
la verificación formal; y el capítulo~\ref{chap:pruebas} recoge los
resultados experimentales.
